\section{Local obstructions for the upper bound}\label{u:appendix}

This appendix supplies the mixed-axis and corner-bit arguments used
in Section~\ref{u:section}. Coordinates, masks, and next-state bits
have the notation of that section. Every pair of starts below is
in state zero. All displayed geometries are induced paths: for an
$L$ or $U$ this follows from its defining coordinates, and for each
listed endpoint order no nonconsecutive pair is at unit distance.

\subsection{Excluding one reset axis}\label{u:appendix-mixed}

\begin{lemma}\label{u:mixed0}
A path survivor cannot have vertical straight map $C_0$ and
horizontal straight map $I$ in the frame of Lemma~\ref{u:axis}.
\end{lemma}

\begin{proof}
Here
$\delta(0,\mathrm{NS})=(\Sdir,0)$,
$\delta(1,\mathrm{NS})=(\N,0)$, and $n_0=e_0=1$.
We force three corner rows and then contradict survival. Every
unlisted corner or endpoint bit remains arbitrary.

\emph{The ES state-zero row must be $(\Sdir,0)$.}
If it chooses east, use $L(\Sdir,\E;3,2)$ with starts $S_2,C$.
The southern agent stays in the reset pocket $S_1,S_2,S_3$.
The corner-started agent stays in $C,E_1,E_2$: every return to $C$
from the horizontal straight cell has state zero, and that corner
row sends it east. The disjoint regions and even-colored starts
give a trap.
If instead the row is $(\Sdir,1)$, use $L(\Sdir,\E;4,1)$
with starts $S_4,C$. The southern agent stays in $S_2,S_3,S_4$,
with $S_2$ visited only in state zero: from $S_3$ either state
reaches $S_2$ or $S_4$ in state zero, and the endpoint returns
to $S_3$ in either state. The other agent has exact cycle
$C_0\to(S_1)_1\to C_0$. The regions are disjoint and the starts
are even. Only $(\Sdir,0)$ remains.

\emph{The SW state-zero row must point west.}
Use the upward U with endpoint order
\[
(0,0),(0,1),(0,2),(1,2),(2,2),(2,1),(2,0)
\]
and starts $(0,1),(2,1)$. Every arrival at either upper corner
from its straight vertical arm has state zero. The ES corner
returns south by the row just forced. If the SW state-zero row
also chose south, both vertical columns would be closed and
two units apart. Thus SW in state zero chooses west.

\emph{The SW state-one row must be $(\Sdir,0)$.}
If it chooses west, use $L(\Sdir,\W;3,2)$ with starts $S_3,W_1$.
The pocket's closed tagged set contains $(S_3)_0$.
The other agent begins
\[
(W_1)_0\to(W_2)_0\to(W_1)_1\to C_1,
\]
using $e_0=1$. Every later arrival at $C$ from the west arm also
has state one, so that arm, including $C$, is closed. The regions
are disjoint and both starts are odd.
If the row is $(\Sdir,1)$, use $L(\Sdir,\W;4,2)$ with starts
$S_4,W_2$. The southern agent remains in $S_2,S_3,S_4$ as above.
The other begins $(W_2)_0\to(W_1)_1\to C_1$ and then
\[
C_1\to(S_1)_1\to C_0.
\]
The already forced west-pointing state-zero row returns it to the
west arm; every return from that arm reaches $C_1$.
Thus its whole trajectory lies in $W_2,W_1,C,S_1$, with $S_1$
visited only in state one. These regions are disjoint and the
starts are even. This excludes the second alternative and forces
$(\Sdir,0)$.

\emph{The forced rows give a trap.}
Use
\[
(0,0),(0,1),(0,2),(1,2),(2,2),(2,1)
\]
with starts $(0,1),(2,1)$. The left three-cell column is closed
by the vertical reset and the ES state-zero row. The right agent
has exact cycle
\[
(2,1)_0\to(2,2)_1\to(2,1)_0
\]
by $n_0=1$ and the SW state-one row. Their horizontal coordinates
remain zero and two, respectively. This contradicts path survival.
\end{proof}

\begin{lemma}\label{u:mixed1}
A path survivor cannot have vertical straight map $C_1$ and
horizontal straight map $I$ in the frame of Lemma~\ref{u:axis}.
\end{lemma}

\begin{proof}
Now
$\delta(0,\mathrm{NS})=(\Sdir,1)$,
$\delta(1,\mathrm{NS})=(\N,1)$, and $e_0=1$.
Reset pockets point north. We keep the initial state zero throughout;
in particular, the free bit $s_1$ is treated in both of its cases.

\emph{The NW state-one row must be $(\N,1)$.}
If it chooses west, use $L(\N,\W;3,2)$ with starts $N_2,W_2$.
The north pocket is closed, while the other agent starts
$(W_2)_0\to(W_1)_1\to C_1$. Every later arrival from the west
arm has state one, so its west-pointing corner row closes
$W_2,W_1,C$. The regions are disjoint and the starts are even.
If the row is $(\N,0)$, use $L(\N,\W;4,1)$ with starts $N_3,W_1$.
The first agent stays in $N_2,N_3,N_4$: the closed pocket tagged
set is $\{(N_2)_1,(N_3)_0,(N_3)_1,(N_4)_1\}$.
The second has prefix $(W_1)_0\to C_1$ followed by
$C_1\to(N_1)_0\to C_1$.
The disjoint regions and odd starts again give a trap.
This leaves only $(\N,1)$.

\emph{The NE state-one row must point east.}
Use the downward U, namely $U(\N,\E)$, with starts
$(0,1),(2,1)$. Both arrivals at the lower corners from their
vertical arms have state one. The right NW corner returns north.
If the left NE state-one row also chose north, both columns would
be closed and two units apart. It must therefore choose east.

\emph{The NE state-zero row must be $(\N,1)$.}
If it also chooses east, use $L(\N,\E;3,2)$ with starts $N_2,E_2$.
The north pocket is closed, and both corner states point east,
closing $C,E_1,E_2$. The regions are disjoint and the starts even.
If the row is $(\N,0)$, use $L(\N,\E;4,2)$ with starts $N_4,E_2$.
The first endpoint move reaches $N_3$ in either state. From there,
state zero reaches $(N_2)_1$ and state one reaches $(N_4)_1$;
thereafter $N_2,N_3,N_4$ remains closed, with $N_2$ only in state
one. The second agent stays in $C,E_1,E_2,N_1$. Its only northern
excursion is $C_0\to(N_1)_0\to C_1$, and the state-one corner
row sends it east. The regions are disjoint and the original
starts even. Hence the remaining row is $(\N,1)$.

\emph{Both values of $s_1$ now give a trap.}
For $s_1=0$, use
\[
(0,1),(0,0),(1,0),(2,0),(2,1),(2,2)
\]
with starts $(0,0),(2,2)$. The left agent has exact cycle
$(0,0)_0\to(0,1)_1\to(0,0)_0$.
The right column is closed because every corner arrival from
its straight cell has state one and the NW row returns north.
The columns are two units apart from time zero.
For $s_1=1$, use $U(\N,\E)$ with starts $(0,0),(2,1)$.
The complete entry prefixes are
\[
\begin{array}{l}
(0,0)_0\to(0,1)_1\to(0,2)_1,\\
(2,1)_0\to(2,0)_1\to(2,1)_1\to(2,2)_1.
\end{array}
\]
Each is then followed by its upper edge in state one. Again the
agents remain in their separate columns throughout. These two
cases exhaust the endpoint bit and contradict path survival.
\end{proof}

\subsection{The four corner output bits}\label{u:appendix-bits}

\begin{lemma}\label{u:bit-laws}
Assume identity straight maps, $n_0=e_0=1$, and the direction
conclusions
\[
\operatorname{dir}\delta(0,\mathrm{NE})=\E,\quad
\operatorname{dir}\delta(0,\mathrm{ES})=\Sdir,\quad
\operatorname{dir}\delta(1,\mathrm{ES})=\E,\quad
\operatorname{dir}\delta(1,\mathrm{SW})=\Sdir.
\]
If the controller survives all induced paths through seven cells,
their output states are, respectively, $1,0,1,0$.
\end{lemma}

\begin{proof}
Each of the following six-cell paths excludes one wrong bit,
independently of the other three output-bit conclusions.

\emph{ES in state zero cannot output one.}
Suppose $\delta(0,\mathrm{ES})=(\Sdir,1)$ and use
\[
(1,2),(0,2),(0,1),(0,0),(1,0),(2,0)
\]
with starts $C=(0,0)$ and $U=(0,2)$.
Write $M=(0,1)$, $H=(1,0)$, $E=(2,0)$, and $P=(1,2)$.
The first agent stays in $C,H,E$: every return from $H$ reaches
$C$ in state zero, whose NE row points east.
The second has closed tagged set
\[
\{U_0,U_1,M_1,P_0,P_1\}.
\]
Indeed $U_0\to M_1\to U_1$, the ES state-one row points to $P$,
and endpoint $P$ returns to $U$ in either state. The two spatial
regions are disjoint and the starts even, proving failure.

\emph{ES in state one cannot output zero.}
Suppose $\delta(1,\mathrm{ES})=(\E,0)$ and use
\[
(0,1),(0,2),(1,2),(2,2),(2,1),(2,0)
\]
with starts $P=(0,1)$ and $I=(2,1)$.
Put $C=(0,2)$, $H=(1,2)$, $D=(2,2)$, $B=(2,0)$.
The left prefix is
$P_0\to C_1\to H_0\to C_0$.
At $C_0$ the ES row points south to $P$, whose endpoint return
may reach either state of $C$. Thus $P,C,H$ is closed, with $H$
only in state zero. The right arm $B,I,D$ is closed: every
arrival at $D$ from $I$ has state one and its SW row points south.
The regions are disjoint and both starts odd.

\emph{NE in state zero cannot output zero.}
Suppose $\delta(0,\mathrm{NE})=(\E,0)$ and use
\[
(0,3),(0,2),(1,2),(2,2),(2,1),(2,0)
\]
with starts $(1,2),(2,1)$.
The first agent has exact cycle $(1,2)_0\leftrightarrow(0,2)_0$.
The second stays in the right three-cell column by the
state-one SW direction and the corner-arrival rule.
The spatial regions are disjoint and both starts odd.

\emph{SW in state one cannot output one.}
Suppose $\delta(1,\mathrm{SW})=(\Sdir,1)$ and use
\[
(0,2),(1,2),(1,1),(1,0),(2,0),(3,0)
\]
with starts $(0,2),(2,0)$.
The upper agent's complete prefix and cycle are
\[
(0,2)_0\to(1,2)_1\to(1,1)_1\to(1,2)_1\to\cdots,
\]
using $e_0=1$. The lower agent stays in the horizontal arm
$(1,0),(2,0),(3,0)$: every arrival at its NE corner $(1,0)$
has state zero and is sent east. The regions are disjoint and
both original starts even. This excludes the last wrong bit.
\end{proof}
